\section{Conclusões sobre a viabilidade da arquitetura}

A revisão do MS-Monitor identifica uma arquitetura coerente para
aquisição de medições, execução local de um classificador compacto,
comunicação por HTTP e persistência dos registros. O código separa
a telemetria de vibração dos atributos térmicos e higrométricos,
trata falhas dos sensores e permite inferência local independente
do envio pela rede. A compilação registrada e a comparação
Python--C++ sustentam a viabilidade computacional preliminar.

Essa conclusão é limitada à implementação e às verificações descritas.
Não demonstra desempenho de detecção de fungos, contaminação ou
micotoxinas. O desbalanceamento do teste, a falta de exemplos de
duas classes e o recall nulo da classe moderada impedem atribuir
validade agronômica ao classificador atual. Seu uso apropriado,
nesta etapa, é a demonstração de aquisição, inferência embarcada
e integração de dados em uma atividade acadêmica.

A continuidade do projeto exige caracterização física da montagem,
comparação metrológica dos sensores, medição do consumo e avaliação
do tempo de inferência na ESP32. A investigação da relação com
condições de deterioração dependerá de dados experimentais
representativos e referências independentes. Controle e classificação
de ventilação permanecem fora da implementação analisada.

\section*{Declaração de apoio de inteligência artificial}

Ferramentas de IA generativa foram utilizadas como apoio à revisão
textual e técnica, à organização do documento LaTeX e à elaboração
de diagramas. As afirmações foram confrontadas com o código e as
fontes indicadas, e exemplos simulados foram identificados como tais.
Esse apoio não substitui a revisão dos integrantes, a avaliação
docente ou os ensaios experimentais. A responsabilidade pela versão
submetida permanece com os autores.
